\section{Графический интерфейс}
\label{sec:frontend}

\subsection{Назначение и устройство окна}
Frontend реализован на Qt Widgets. Его назначение --- получить параметры
эксперимента, вызвать backend и показать результаты. Точка входа
\code{main.cpp} создаёт \code{QApplication}, экземпляр \code{MainWindow},
показывает окно и запускает цикл событий. Начальная выбранная вкладка в
\code{mainwindow.ui} --- <<Анимация>>; до настройки и запуска модель ещё
не инициализирована.

Разметка окна хранится в \code{mainwindow.ui}. Инструмент Qt \code{uic}
создаёт из неё \code{ui_mainwindow.h}. В конструкторе вызывается
\code{ui->setupUi(this)}. Логика кнопок, чтение параметров и обновление вкладок
находятся в \code{mainwindow.cpp}. Расчёты склада, цены и прибыли выполняет
backend; таблицы читают полученные данные.

\begin{tabularx}{\textwidth}{@{}L{0.38\textwidth}Y@{}}
\toprule
Модуль & Ответственность \\
\midrule
\code{mainwindow.ui} & Разметка, кнопки, поля параметров, контейнер вкладок. \\
\code{mainwindow.h/.cpp} & Слоты кнопок, таймер дней, передача параметров,
вызовы модели и обновление представлений. \\
\code{front/*tab.h/.cpp} & Реализации отдельных вкладок. \\
\code{resources.qrc}, \code{images/} & Встроенные изображения склада,
дома, курьера и телефона. \\
\code{front/widgets/loadring.*}, \code{front/widgets/stockbar.*}
& Заготовки с пустыми конструкторами; текущие вкладки их не используют. \\
\bottomrule
\end{tabularx}

\subsection{Порядок работы пользователя}
\begin{enumerate}
  \item Открыть <<Параметры>> и задать условия эксперимента.
  \item Нажать <<Сгенерировать>>: модель инициализируется, очищается журнал,
        показываются начальные партии; окно переключается на <<Склад>>.
  \item Для пошаговой работы нажимать <<Шаг>>. Для автоматической --- <<Старт>>.
        При нулевом дне <<Старт>> сам повторно выполняет инициализацию.
  \item Во время расчёта смотреть <<Склад>>, <<Заказы>>, <<Курьеры>>,
        <<Заявки>>, <<Журнал>> и <<Анимация>>. <<Пауза>> останавливает таймер дней.
  \item После дня $N$ автоматический таймер останавливается. Итоги доступны
        на вкладке <<Результаты>>.
  \item Для нового эксперимента снова задать параметры и нажать
        <<Сгенерировать>>. Это заново создаёт состояние backend.
\end{enumerate}
До первой инициализации нельзя начинать работу кнопкой <<Шаг>>: она сама
не вызывает \code{init_simulation()}. Для одиночных шагов во время
автоматического запуска сначала используйте <<Пауза>>.

\subsection{Кнопки и элементы управления}
\begin{longtable}{@{}L{0.27\textwidth}L{0.66\textwidth}@{}}
\toprule
Элемент & Фактическое поведение \\
\midrule
\endhead
<<Сгенерировать>> & Останавливает таймер, ставит день 0, читает форму,
вызывает \code{configure_simulation()} и \code{init_simulation()},
обновляет вкладки и сбрасывает анимацию. \\
<<Старт>> & При дне 0 вызывает генерацию; затем запускает таймер дней.
При положительном дне продолжает отсчёт с текущим состоянием. \\
<<Пауза>> & Останавливает таймер \code{MainWindow}; backend не выполняет
следующие автоматические шаги. Отдельный таймер анимации не останавливается. \\
<<Шаг>> & Непосредственно вызывает \code{doStep()} для следующего дня;
сам не инициализирует модель и не выключает автоматический таймер. \\
<<Сброс>> & Останавливает таймер, сбрасывает отображаемый день и журнал.
Сам по себе не очищает состояние backend; для нового эксперимента
нужна последующая генерация или старт с дня 0. \\
Поле <<День>> & Показывает счётчик \code{MainWindow::day}. Ручное изменение
текста не изменяет внутренний номер дня модели. \\
Ползунок <<Скорость>> & Меняет реальный интервал между дневными шагами,
но не длительность модельного дня и не скорость отдельной анимации. \\
\bottomrule
\end{longtable}

\subsection{Параметры и связь с backend}
\code{MainWindow::buildParameters()} переносит поля формы в
\code{SimulationParameters}. Названия, диапазоны и отличия формы от
backend приведены в разделе~\ref{sec:running}. В частности,
\code{ordersPerDaySpin} по умолчанию равен 20, а вероятность из
\code{discountedProbSpin} делится на 100. Число постоянных покупателей
в форме задано кодом как 10, seed --- 5489.
Изменения настроек применяются при следующей инициализации; уже настроенная
модель не перечитывает поля на каждом шаге.

Кнопки соединяются со слотами по соглашению Qt:
\code{on_<objectName>_clicked()}, через настройку формы.
\code{QTimer::timeout} явно соединён с \code{MainWindow::tick()}, а
\code{QSlider::valueChanged} --- с \code{onSpeedChanged()}.
Один таймерный вызов соответствует одному дню:
\begin{lstlisting}[language=C++]
manage_warehouse(day);
generate_daily_events(day);
process_daily_order(day);
cleanup_random_customers();
refreshAllTabs();
\end{lstlisting}
В реальном \code{doStep()} между завершением backend и обновлением вкладок
также добавляется строка журнала. Номер дня увеличивается перед вызовами
backend. Исключения при инициализации и шаге перехватываются и записываются
в журнал; при ошибке шага таблицы не обновляются. Для восстановления
согласованного состояния следует остановить таймер и выполнить генерацию.

Для значения ползунка $v=1,\ldots,10$ интервал таймера:
\begin{equation}
  \Delta t=\max(100,2000-210(v-1))\text{ мс}.
\end{equation}
На используемом диапазоне это 2000--110 мс; начальное значение $v=5$
даёт 1160 мс. Таймер выполняет шаги в потоке интерфейса; отдельного
вычислительного потока или сетевого обмена frontend/backend нет.

\subsection{Содержимое вкладок}
\begin{longtable}{@{}L{0.22\textwidth}L{0.71\textwidth}@{}}
\toprule
Вкладка & Источник и состав данных \\
\midrule
\endhead
<<Параметры>> & Значения настроек и кнопка генерации. \\
<<Склад>> & \code{warehouse.inventory}, каталог и справочники из
\code{back/data.h}. Одна строка --- партия: ID, форма лекарства, группа,
дозировка, количество, день истечения срока, текущая цена и оптовая цена. \\
<<Заказы>> & \code{simulation_result().completed_orders}. Только доставленные
покупки: номер, день заказа, день доставки, курьер, ФИО, адрес, запрошенные
и выданные количества, цена. Очередь и подготовленные покупки отдельными
строками здесь не отображаются. \\
<<Курьеры>> & Дневная нагрузка последнего завершённого дня, сумма по дням
и признак попадания в 7--15. Нумерация на экране начинается с 1. \\
<<Заявки>> & \code{pending_supplies}: ID лекарства, название с дозировкой,
количество и оставшееся число дней до поставки. \\
<<Результаты>> & Текстовый отчёт: доход, расходы, списания, прибыль,
стоимость остатков, количества начальных, поступивших и списанных упаковок,
число доставок и размеры ожидающих очередей. \\
<<Журнал>> & Текстовые сообщения инициализации, шагов и ошибок. \\
<<Анимация>> & Условная иллюстрация склада, города и движения курьеров;
текущий день, суммарный складской остаток и индикатор новых заявок. \\
\bottomrule
\end{longtable}

Таблицы построены на \code{QTableView} и \code{QStandardItemModel};
их редактирование выключено. Метод \code{refresh()} удаляет старые строки,
затем заново заполняет их. Текстовые вкладки используют \code{QTextEdit}
в режиме чтения. \code{refreshAllTabs()} вызывается после инициализации
и успешного дня; фактические расчётные значения читаются из backend.

\subsection{Устройство анимации}
В \code{AnimTab} работает отдельный таймер с интервалом 16 мс,
соединённый с \code{animate()}. Метод \code{paintEvent()} рисует сцену
через \code{QPainter}. Изображения \code{QPixmap} загружаются из ресурсов
\code{:/images/warehouse.png}, \code{house.png}, \code{courier.png}
и \code{phone.png}; при отсутствии используются простые фигуры.

После завершённого дня \code{refreshAllTabs()} передаёт
\code{updateState(day, stockCount, courier_load, created_supplies)}.
При положительной нагрузке неподвижный курьер начинает поездку от склада
к городу, возвращается и снова останавливается. Его состояния:
\code{Idle}, \code{GoingToCity}, \code{Returning}.
Координата $x$ меняется на условную скорость 3--5 пикселей за такт,
а $y$ получает синусоидальное отклонение. При новых заявках поставщику
на 120 тактов включается мигающий значок телефона.

Вложенная структура \code{AnimTab::Courier} хранит только визуальные
координаты, состояние движения и нагрузку. Она не является курьером
расчётной модели: в backend курьеры по-прежнему представлены числом $M$
и индексами массива. Движение не рассчитывает адреса, длительность поездок
и стоимость доставки и не изменяет результаты backend.

\subsection{Особенности текущей версии}
Документация описывает восстановленный интерфейс, а не ранее добавленную
городскую карту. В нём нет выбора количества зданий, масштаба карты,
кнопки экспорта JSON, загрузки CSV и редактирования seed.
JSON и CSV доступны через программный интерфейс backend при отдельном
подключении соответствующих возможностей.

При проверке экрана учитывайте следующие особенности реализации:
\begin{itemize}
  \item <<Сброс>> не равен полной инициализации backend; старые результаты
        могут оставаться видимыми до следующей генерации.
  \item Пауза моделирования не останавливает уже начатое визуальное движение.
  \item Столбец <<Цена партии>> на складе фактически показывает
        \code{batch_price}, то есть цену \emph{упаковки}, а не стоимость всей партии.
  \item Надпись <<Партий изначально>> в результатах выводит
        \code{initial_packages}, то есть количество \emph{упаковок}.
  \item В \code{RequestsTab::refresh()} ID лекарства сравнивается с длиной
        справочника форм вместо каталога; поэтому для части заявок название
        отображается как <<?>>. Количество и таймер берутся из заявки.
  \item В формате дневной строки журнала повторён плейсхолдер
        \texttt{\%1}; некоторые счётчики строки поэтому не соответствуют
        переданным аргументам. Для точных значений следует смотреть
        \code{DailySimulationResult}, таблицы и \code{SimulationResult}.
\end{itemize}
Эти особенности представления не меняют формулы backend. Их исправление
требует отдельного изменения программы; добавление этого раздела
не изменяет исходники интерфейса.
